\chapter{La naturaleza}
\label{ch:nature}

\section{Un accidente que vale la pena conservar}

Cuba tiene la ecología insular grande menos degradada del Caribe, y la tiene en buena medida por
accidente.

El país conserva extensos sistemas de arrecife de coral, entre ellos los Jardines de la Reina
---un archipiélago frente a la costa sur, cerrado a la pesca comercial desde 1996, donde las
poblaciones de tiburones y meros están entre las últimas de la región que se parecen a lo que
eran los arrecifes caribeños antes de la pesca industrial---. Conserva la Ciénaga de Zapata, el
mayor humedal del Caribe insular, con una lista de aves que incluye varias endémicas que no
existen en ninguna otra parte. Conserva el Parque Nacional Alejandro de Humboldt en el oriente,
uno de los sitios insulares biológicamente más diversos del planeta por superficie, y los
mogotes de Viñales. El endemismo cubano es alto en plantas, aves, reptiles y moluscos, como
ocurre en cualquier isla grande y antigua, y una parte inusualmente grande de él sigue ahí.

Las razones no son virtuosas. Cuba dejó de poder pagar fertilizante y plaguicida en 1990. Dejó de
poder construir hoteles costeros a los ritmos caribeños. Nunca desarrolló el litoral de cruceros
y condominios que ha consumido buena parte de la región. Sus ciudadanos no podían comprar botes
ni segundas viviendas. La pobreza es una política de conservación eficaz, brutal y temporal, y
cada año de la recuperación que este libro propone la volverá menos eficaz.

Hay también un elemento institucional real, que hay que reconocer. Cuba protegió deliberadamente
los Jardines de la Reina y lo hizo cumplir. Tiene un sistema nacional de áreas protegidas que
cubre una parte sustancial del territorio, una ley ambiental seria de 1997, una tradición genuina
de biología de campo, y en la Tarea Vida un programa nacional de adaptación climática adoptado en
2017 que es científicamente mejor que la mayoría de los de la región y está casi enteramente sin
financiar. El conocimiento existe. El dinero y la capacidad de hacer cumplir, no.

\section{El mecanismo que lo destruye}

El patrón es el mismo en todo el Caribe y vale la pena enunciarlo como mecanismo y no como
advertencia, porque el mecanismo es lo que un diseño puede interrumpir.

Un país decide hacer crecer el turismo. El frente de playa se convierte en la tierra más valiosa,
así que se construye sobre él ---y construir sobre él elimina la duna, el manglar y el pasto
marino, que era lo que sujetaba la arena y filtraba la escorrentía---. La construcción y la nueva
población producen aguas residuales que la infraestructura de tratamiento no maneja, así que
entran nutrientes al agua. Los nutrientes alimentan algas; las algas ahogan el coral. La pesca se
intensifica para alimentar a los visitantes, eliminando los peces herbívoros que pastan las algas
sobre el coral. Los botes anclan en el arrecife. Las carreteras y los pedraplenes cortan la
circulación del agua en las lagunas detrás de los cayos. En quince a veinticinco años el arrecife
con el que se vendía el destino está muerto, la playa que protegía se está erosionando, y el país
defiende su litoral con roca importada.

Cuba está al principio de esta secuencia y no al final, y tiene una ventaja enorme: puede ver la
curva entera, en detalle, en una docena de jurisdicciones a una hora de vuelo.

\section{La exposición}

El riesgo climático en Cuba no es un escenario futuro; es un calendario.

Los \emph{huracanes} son el riesgo agudo, y la tendencia en la intensidad de las tormentas más
fuertes es al alza. La defensa civil cubana es genuinamente excelente y sus cifras de muertos
están entre las más bajas de la región para tormentas de fuerza comparable ---es un logro cubano
real y poco apreciado--- pero el daño físico no se previene evacuando. La temporada de 2008 por
sí sola costó del orden de una quinta parte del producto nacional.

El \emph{nivel del mar} es el riesgo crónico. Las evaluaciones de la propia Tarea Vida cubana
proyectan la pérdida de asentamientos costeros bajos a lo largo del siglo e identifican
comunidades que no deberían reconstruirse donde están ---cosa políticamente extraordinaria de
publicar para un gobierno, y que publicó---.

La \emph{intrusión salina} en los acuíferos costeros ya es medible y es un problema de
abastecimiento de agua para La Habana y para la llanura azucarera.

El \emph{blanqueamiento del coral} sigue a la temperatura del mar, y el Caribe ha tenido eventos
regionales repetidos. Los arrecifes cubanos no están exentos; lo que tienen es mejor condición de
partida, que es el determinante individual más importante de si un arrecife se recupera de un
evento de blanqueamiento.

La \emph{sequía} en las provincias orientales ---las más pobres, según el
capítulo~\ref{ch:premises}--- ha sido severa y recurrente, e interactúa con todo lo del
capítulo~\ref{ch:land}.

\section{El diseño}

El argumento en todo momento es económico y no sentimental, porque la versión sentimental pierde
frente a una propuesta hotelera siempre.

Los arrecifes, los humedales y los cayos son el acervo físico desde el cual se venden los
productos de alto valor del capítulo~\ref{ch:tourism}. Un día de buceo en los Jardines de la Reina
gana más por visitante que cualquier habitación de playa del país y existe únicamente porque
alguien trazó una línea y la hizo cumplir. Eso no es una restricción sobre la estrategia
turística; es el inventario de la estrategia turística, y el tratamiento correcto del inventario
no es consumirlo.

\begin{design}{Proteger el acervo}
\label{d:protect}
\begin{rules}
\rl{1}{Un regulador independiente} Una autoridad ambiental con el mandato, el presupuesto y las
garantías de publicación del Diseño~\ref{d:courts}, independiente del ministerio de turismo, de
las entidades del Diseño~\ref{d:gaesa}, y de cualquier ministerio que sea dueño de un desarrollo.
Las decisiones de licencia y sus fundamentos publicados; las denegaciones publicadas. Hoy el
promotor y la autoridad que autoriza son con frecuencia el mismo Estado.

\rl{2}{Retiro costero, sin excepción} Un retiro legal de construcción respecto de la línea media
de pleamar, aplicado a todo promotor incluidos los estatales y los vinculados a los militares, sin
procedimiento de dispensa. Toda regla de retiro costero del Caribe que tiene procedimiento de
dispensa ha fracasado a través de él, y por eso esta cláusula trata sobre la ausencia de la
excepción antes que sobre la distancia.

\rl{3}{Ningún pedraplén nuevo} No se construyen más pedraplenes hacia los cayos; los existentes se
evalúan para su remediación con alcantarillas que restauren la circulación. El daño que han hecho
a los sistemas lagunares está documentado y es en parte reversible.

\rl{4}{Aguas residuales antes que habitaciones} Ningún desarrollo de alojamiento se autoriza sin
capacidad de tratamiento puesta en servicio y operando. La carga de nutrientes es el mayor
impulsor controlable de la pérdida de coral, y es una decisión de infraestructura que se toma años
antes de que muera el arrecife.

\rl{5}{La pesca} Zonas de veda total mantenidas y ampliadas en torno a los sistemas de arrecife,
con la vigilancia financiada con los ingresos de visitantes según la cláusula~1 del
Diseño~\ref{d:natfinance}. El modelo de los Jardines de la Reina funciona y es replicable.

\rl{6}{Datos ambientales publicados} Estado del arrecife, calidad del agua, extensión de las áreas
protegidas y acciones de vigilancia publicados anualmente, según el Diseño~\ref{d:minimum}.3.
\end{rules}
\end{design}

\begin{design}{Cómo se paga}
\label{d:natfinance}
\begin{rules}
\rl{1}{Financiamiento por visitantes} Una proporción del gravamen al alojamiento y de las tasas de
acceso a áreas protegidas se asigna por ley al sistema de áreas protegidas y la retiene este, al
modo del Diseño~\ref{d:heritage}. La conservación que depende de una negociación presupuestaria
anual pierde la negociación.

\rl{2}{Ponerle precio al acceso} Donde un sitio esté en su capacidad de carga, el acceso se cobra y
se asigna por subasta o sorteo publicados (Diseño~\ref{d:limits}). Esto convierte la escasez en
ingreso para aquello que es escaso.

\rl{3}{Carbono y servicios ecosistémicos} El carbono de manglares y bosques, y los esquemas de pago
por servicios ecosistémicos, perseguidos como fuente de ingresos. Son mercados imperfectos y los
montos son modestos, pero son divisa por no hacer algo, lo que se ajusta inusualmente bien a la
posición de Cuba.

\rl{4}{La ciencia como exportación} La biología de campo cubana más los ecosistemas intactos más
el financiamiento extranjero a la investigación es una línea de negocio real e inmediata, y se
solapa con el visitante de alto valor del capítulo~\ref{ch:tourism} y con la apertura académica del
capítulo~\ref{ch:talent}. Solo exige que los científicos extranjeros puedan obtener permisos sin
una evaluación política.
\end{rules}
\end{design}

\begin{design}{Construir para el siglo}
\label{d:adapt}
\begin{rules}
\rl{1}{Financiar la Tarea Vida} El programa nacional de adaptación existente es científicamente
sólido y no está financiado. Se le asigna una partida permanente del fondo de estabilización del
Diseño~\ref{d:rentrule}.2, que es exactamente el tipo de gasto de capital para el que existe ese
fondo.

\rl{2}{Códigos de edificación} Normas de construcción resistente a huracanes, aplicadas a través
del control municipal de la edificación del Diseño~\ref{d:municipal}.3, y aplicadas a la vivienda
tanto como a los hoteles. El parque habitacional cubano es su activo más vulnerable a las
tormentas, y su reconstrucción bajo el capítulo~\ref{ch:capital} es la oportunidad de una vez por
siglo de elevar el estándar.

\rl{3}{Retirada gestionada} Donde la Tarea Vida identifique asentamientos que no pueden
defenderse, la reubicación se planifica, se financia, se consulta y se compensa, y no se financia
públicamente ninguna reconstrucción en el sitio. Es políticamente doloroso y más barato que la
alternativa, y Cuba ya ha hecho el trabajo analítico.

\rl{4}{El agua} Protección de acuíferos, reducción de fugas en la red de distribución ---donde las
pérdidas son muy grandes--- y precio del agua por encima de una asignación básica gratuita. Cuba
pierde hoy una gran cantidad de agua bombeada en fugas, lo que es a la vez un problema de agua y,
como el agua se bombea, un problema de electricidad según el capítulo~\ref{ch:energy}.

\rl{5}{Seguro y reserva} Según el Diseño~\ref{d:storm}.5, una reserva explícita y publicada contra
pérdidas catastróficas, puesto que la cobertura comercial para activos cubanos es cara y en varios
mercados inobtenible.
\end{rules}
\end{design}

\section{La disyuntiva, enunciada}

Este capítulo restringe al capítulo~\ref{ch:tourism} y al capítulo~\ref{ch:capital} y debería
decirlo en vez de fingir que los objetivos se alinean automáticamente.

Los retiros costeros, los requisitos de tratamiento de aguas, los límites de capacidad de carga y
las zonas de veda reducen el número de habitaciones de hotel que pueden construirse y la velocidad
a la que pueden construirse. A corto plazo eso es menos inversión, menos empleo en construcción y
menos ingresos, y a un gobierno bajo presión fiscal le ofrecerán mucho dinero por hacer una
excepción ---por lo cual el Diseño~\ref{d:protect}.2 no tiene procedimiento de dispensa y por lo
cual la independencia del regulador está escrita en los mismos términos que la de la judicatura---.

El argumento para aceptar la restricción es el del capítulo~\ref{ch:tourism}: Cuba no compite en
volumen, no puede ganar en volumen, y se propone vender lo único genuinamente escaso que hay en el
Caribe. Todos los demás países de la región ya hicieron el canje contrario y pueden informar del
resultado.
